% =====================================================================
% cap-01b-trabalhos-relacionados.tex — Capítulo 3: Trabalhos Relacionados
% =====================================================================

\chapter{Trabalhos Relacionados}
\label{cap:relacionados}

Este capítulo situa o presente trabalho no panorama de sistemas de
caixão de areia com realidade aumentada, revisando a origem do
conceito e sua adoção acadêmica (Seção~\ref{sec:estado-arte}) e
posicionando explicitamente as contribuições deste PFC frente a esse
estado da arte (Seção~\ref{sec:posicionamento}). A diversificação de
sensores de profundidade empregados por essas implementações já foi
discutida na Seção~\ref{subsec:sensores-alt}.

\section{AR Sandboxes: Estado da Arte e Adoção Acadêmica}
\label{sec:estado-arte}

O conceito de caixão de areia com realidade aumentada foi consolidado
pelo projeto \textbf{SARndbox}, desenvolvido por
Kreylos~\citeonline{sarndbox} no KeckCAVES da Universidade da Califórnia em
Davis. Escrito em C\texttt{++}/OpenGL para aproveitar processamento
gráfico dedicado, o SARndbox captura a superfície da areia com um
Kinect, projeta escalas hipsométricas e curvas de nível coloridas em
tempo real e simula, sobre o relevo físico, o escoamento de água por
um solver de dinâmica de fluidos. Mais de 150 instalações seguindo
esse modelo foram documentadas em escolas, universidades, museus e
centros de pesquisa, e a própria página do projeto lista, entre os
domínios de aplicação já explorados pela comunidade, hidrologia,
geologia, mineração, planejamento regional e \textbf{operações
militares}~\cite{sarndbox} — ainda que, na revisão realizada para
este trabalho, não tenham sido encontradas implementações
publicamente documentadas que combinem esse último domínio com
comparação quantitativa a um Modelo Digital de Elevação cartográfico
oficial, lacuna que motiva a Seção~\ref{sec:posicionamento}.

A validade pedagógica da abordagem foi avaliada empiricamente por
Woods et al.~\citeonline{woods2016pilot}, em um estudo piloto que empregou o
SARndbox em laboratórios introdutórios de Geologia para ensinar
leitura de mapas topográficos e processos superficiais, relatando
ganhos observados na capacidade de visualização espacial dos
estudantes. Esse resultado é relevante para o presente trabalho por
oferecer evidência prévia, independente e publicada, de que a
transição cognitiva entre carta topográfica e relevo tridimensional
tangível — o problema central motivador deste PFC
(Seção~\ref{sec:contexto}) — é de fato facilitada pela metáfora do
caixão de areia aumentado, dispensando a necessidade de este trabalho
revalidar esse pressuposto pedagógico isoladamente.

Mais recentemente, Wellmann et al.~\citeonline{wellmann2022openarsandbox} propuseram o
\textbf{Open~AR-Sandbox}, uma reimplementação de código aberto em
Python voltada à visualização de informação geológica de subsuperfície
(camadas de litologia, não apenas topografia), com interface
modularizável para uso em pesquisa e divulgação científica. A escolha
por Python, em vez do C\texttt{++} nativo do SARndbox, aproxima essa
implementação do ecossistema científico também adotado neste trabalho
(NumPy, OpenCV — Seção~\ref{sec:tradeoffs}), reforçando que a
linguagem não é mais um fator limitante de desempenho para esse tipo
de aplicação em hardware convencional.

\section{Posicionamento do Presente Trabalho}
\label{sec:posicionamento}

A Tabela~\ref{tab:relacionados} resume as diferenças entre as
implementações revisadas e o sistema desenvolvido neste PFC, ao longo
de quatro eixos de comparação.

\begin{table}[!ht]
\centering
\caption{Posicionamento do presente trabalho frente a implementações revisadas.}
\label{tab:relacionados}
\small
\begin{tabularx}{\textwidth}{@{}l X X@{}}
\toprule
\textbf{Eixo} & \textbf{SARndbox e variantes revisadas} & \textbf{Este trabalho} \\
\midrule
\textit{Feedback} ao operador &
Curvas de nível, escalas hipsométricas ou camadas geológicas —
qualitativo, sem alvo externo definido &
Decisão célula a célula (cavar/preencher/manter) por comparação
quantitativa com tolerância explícita contra um MDE de referência \\
\addlinespace
Resiliência a falta de hardware &
Não documentada nas publicações revisadas; presume-se sensor e GPU
dedicada disponíveis &
Cadeia de \textit{fallback} em quatro níveis
(Seção~\ref{subsec:fallback}), com simulação interativa completa na
ausência de qualquer sensor \\
\addlinespace
Origem do MDE de referência &
Nenhuma (SARndbox, Woods et al.) ou dado geológico de subsuperfície
sem vínculo cartográfico oficial (Open~AR-Sandbox) &
Arquivo GeoTIFF georreferenciado fornecido pelo operador — compatível
com produtos oficiais como os do BDGEx (Seção~\ref{sec:mde}) —, com
\textit{fallback} sintético para demonstração sem conectividade \\
\addlinespace
Domínio de aplicação &
Educação em geociências, hidrologia e divulgação científica &
Treinamento militar em leitura de carta e planejamento tático, em
atendimento a uma demanda real da Seção de Simulação da AMAN \\
\bottomrule
\end{tabularx}
\end{table}

A contribuição deste trabalho não está, portanto, em propor a técnica
de caixão de areia com realidade aumentada — consolidada desde o
SARndbox e validada pedagogicamente por Woods et al.~\citeonline{woods2016pilot} —,
mas em preencher as três lacunas já enunciadas na
Seção~\ref{sec:contexto}: substituir o retorno visual qualitativo por
uma regra de negócio quantitativa contra um MDE cartográfico oficial,
tornar explícita a decisão operacional célula a célula, e assegurar
resiliência operacional total diante da ausência de hardware — nenhuma
das quais documentada em conjunto nas implementações revisadas neste
capítulo.
